\section{Objet et portée de l'annexe}
\label{annexe-K-objet}

Cette annexe formalise l'équivalence de présentation de la plasticité.
Elle répond à une difficulté de principe : une modification apparente de
la loi d'action peut toujours, dans une classe causale et markovienne
déclarée, être représentée comme l'évolution d'une variable d'état
supplémentaire lue par une loi fixe. L'observation des exécutions peut
alors détecter qu'un résumé d'état est incomplet, mais elle ne peut pas
trancher l'emplacement ontologique de la variable manquante.

Le résultat est relatif à une classe d'histoires $\mathfrak B$, à une
résolution $\rho$ et à un domaine réalisable. Il ne dit pas que toute
mémoire est de dimension finie, ni que toute dépendance historique est
markovianisable. Sa frontière exacte est la réalisabilité causale et
markovienne de la plasticité considérée.

\begin{encadre}[title={Résultat directeur}]
Soit une présentation plastique dans laquelle une variable interne $\gamma$ modifie la loi d'action. Si $\gamma$ possède une dynamique causale, markovienne et réalisable, alors cette présentation est équivalente, pour toute observation fondée sur les exécutions, à une
présentation où $\gamma$ est promue coordonnée d'état et où la loi d'action est figée. Les deux présentations engendrent les mêmes exécutions ; elles sont donc équivalentes en réponse pour \((\mathfrak B,\rho)\).
\end{encadre}

\section{Deux présentations d'un même mécanisme}
\label{annexe-K-deux-presentations}

\subsection{Présentation plastique}

On considère un état de base $x\in X$ et une variable interne
$\gamma\in\Gamma$. Dans la présentation dite \emph{plastique}, la loi
d'action dépend de $\gamma$ :
\[
\dot x=F_{\gamma}(x,u),
\qquad
x^+=J_{\gamma,e}(x),
\]
où $u$ est l'entrée causale et $e$ le type d'événement. Le typage des
événements --- notamment chargé ou vide dans l'instance canonique --- est
lui-même une fonction déclarée de l'état et de l'entrée, identique dans
les deux présentations. La variable $\gamma$ évolue selon
\[
\dot\gamma=G(\gamma,x,u),
\qquad
\gamma^+=K_e(\gamma,x).
\]
L'expression ``la loi change'' désigne ici la convention selon laquelle
$\gamma$ est rangée du côté de la loi plutôt que dans l'état déclaré.

\subsection{Présentation à loi figée}

On définit l'état augmenté
\[
\widetilde x=(x,\gamma)\in \widetilde X:=X\times\Gamma.
\]
La loi figée sur $\widetilde X$ est
\[
\widetilde F((x,\gamma),u)
=
\bigl(F_{\gamma}(x,u),G(\gamma,x,u)\bigr),
\]
et, pour chaque événement $e$,
\[
\widetilde J_e(x,\gamma)
=
\bigl(J_{\gamma,e}(x),K_e(\gamma,x)\bigr).
\]
Le symbole de la loi ne varie plus : toute la variation est portée par
la coordonnée d'état $\gamma$.

\begin{encadre}[title={Convention de présentation}]
Les termes \emph{loi plastique} et \emph{loi figée à état augmenté} ne
désignent pas deux dynamiques différentes. Ils désignent deux découpages
du même mécanisme entre ``état'' et ``loi''. L'équivalence ci-dessous
porte précisément sur ce rebaptême.
\end{encadre}

\section{Hypothèses de l'équivalence}
\label{annexe-K-hypotheses}

Soit $D\subseteq X\times\Gamma$ le domaine réalisable des conditions
initiales et paramètres de protocole. On suppose :

\begin{description}
\item[(P1) Causalité.] Les flots, événements et mises à jour de $\gamma$
ne dépendent que de l'état courant, de l'entrée passée et présente et
des paramètres déclarés avant l'exécution.

\item[(P2) Fermeture markovienne.] Le couple $(x,\gamma)$ détermine
l'évolution future sous tout protocole admissible. Aucune information
historique supplémentaire n'est requise.

\item[(P3) Réalisabilité.] Les états et transitions considérés sont
atteignables dans le domaine déclaré ; les deux présentations utilisent
le même ensemble initial et les mêmes protocoles.

\item[(P4) Déterminisme ou couplage commun.] Pour une condition initiale
et une histoire données, les deux présentations définissent une même
exécution déterministe. Dans une extension stochastique, elles doivent
être couplées par la même réalisation du bruit.

\item[(P5) Observation par les exécutions.] La lecture
$\omega_\rho$ et la réponse $\mathcal R^{\mathfrak B,\rho}$ dépendent de
l'exécution produite, non de l'étiquette ``état'' ou ``loi'' attachée à
ses coordonnées.

\item[(P6) Finitude locale des événements.] Toute exécution admissible ne
contient qu'un nombre fini d'événements sur chaque intervalle de temps
compact ; les accumulations de type Zeno sont exclues du domaine du
théorème.
\end{description}

(P5) est l'axiome d'observation déjà utilisé dans les Appendices G et H.
Les hypothèses nouvelles sont surtout (P1)--(P3). Elles isolent la limite
exacte du résultat.

\section{Correspondance trajectorielle}
\label{annexe-K-correspondance}

Pour une histoire admissible $\mathcal H\in\mathfrak B$, notons
\[
\operatorname{Exec}_{\mathrm{pl}}(x_0,\gamma_0;\mathcal H)
\]
l'exécution de la présentation plastique et
\[
\operatorname{Exec}_{\mathrm{fig}}((x_0,\gamma_0);\mathcal H)
\]
celle de la présentation à loi figée.

\resultstep{lem:K-identite-executions}
\begin{encadre}[title={Lemme \theproposition{} --- Identité des exécutions}]
Sous (P1)--(P4) et (P6), pour toute condition initiale réalisable et toute
histoire admissible,
\[
\operatorname{Exec}_{\mathrm{pl}}(x_0,\gamma_0;\mathcal H)
=
\operatorname{Exec}_{\mathrm{fig}}((x_0,\gamma_0);\mathcal H),
\]
après l'identification canonique de la variable interne $\gamma$ avec la
coordonnée supplémentaire de l'état augmenté.
\end{encadre}

\paragraph{Démonstration.}
Entre deux événements, les deux présentations satisfont par construction
le même système couplé
\[
(\dot x,\dot\gamma)
=
\bigl(F_{\gamma}(x,u),G(\gamma,x,u)\bigr)
\]
avec la même condition initiale et la même entrée. L'unicité locale donne
la même portion de trajectoire jusqu'au prochain événement. Le prédicat
d'événement étant évalué sur les mêmes variables, le prochain événement
survient au même instant et porte le même type. Les deux mises à jour
sont alors identiques :
\[
(x^+,\gamma^+)
=
\bigl(J_{\gamma,e}(x),K_e(\gamma,x)\bigr).
\]
Par (P6), une induction sur les événements rencontrés sur tout intervalle
compact donne l'identité de l'exécution sur cet intervalle, puis sur
l'exécution entière par exhaustion du temps. \hfill$\square$

\section{Théorème d'équivalence de présentation}
\label{annexe-K-theoreme}

\resultstep{thm:K-equivalence-plasticite}
\begin{encadre}[title={Théorème \theproposition{} --- Équivalence de présentation de la plasticité}]
Sous (P1)--(P6), la présentation plastique
\[
(\Psi,\mathcal G_{\mathrm{pl}})
\]
et la présentation à état augmenté et loi figée
\[
(\Psi\oplus\gamma,\mathcal G_{\mathrm{fig}})
\]
sont $\mathcal R$-équivalentes pour toute classe d'histoires
$\mathfrak B$ et toute résolution $\rho$ compatibles avec le domaine
déclaré :
\[
(\Psi,\mathcal G_{\mathrm{pl}})
\sim_{\mathfrak B,\rho}
(\Psi\oplus\gamma,\mathcal G_{\mathrm{fig}}).
\]
Plus précisément, pour toute $\mathcal H\in\mathfrak B$,
\[
\mathcal R_{\mathrm{pl}}^{\rho}(\mathcal H)
=
\mathcal R_{\mathrm{fig}}^{\rho}(\mathcal H).
\]
\end{encadre}

\paragraph{Démonstration.}
Le Lemme~\ref{lem:K-identite-executions} donne l'identité des exécutions pour toute histoire
admissible. Par (P5), la lecture à résolution $\rho$ prend donc la même
valeur dans les deux présentations. L'égalité vaut point par point sur
$\mathfrak B$, donc les cartes de réponse et les quotients d'observation
coïncident. \hfill$\square$

\begin{encadre}[title={Ce que le théorème ne suppose pas}]
Le théorème ne suppose ni que $\gamma$ soit scalaire, ni qu'elle soit
lente, ni que sa dynamique soit linéaire, ni que la sortie soit binaire.
Il ne suppose pas non plus que $\gamma$ soit directement observable. Il
exige seulement qu'elle puisse être incluse dans un état causal et
markovien réalisable et que l'observation porte sur les exécutions.
\end{encadre}

\section{Détection de la plasticité et indécidabilité de son support}
\label{annexe-K-detection}

L'équivalence de présentation n'empêche pas la détection d'une variable
active supplémentaire. Elle interdit seulement de décider si cette
variable ``appartient à la loi'' ou ``appartient à l'état''.

Soit $z_4=(p,\chi,\theta,s)$ le résumé lent de l'Appendice I et
\[
z_5=(p,\chi,\theta,s,\gamma)
\]
son extension plastique. Supposons qu'il existe deux histoires
$\mathcal H_1,\mathcal H_2$ telles que
\[
z_4(\mathcal H_1)=z_4(\mathcal H_2),
\qquad
\gamma(\mathcal H_1)\neq\gamma(\mathcal H_2),
\]
et une sonde $\mathcal P$ pour laquelle
\[
\widehat\Omega_{\mathcal P,\rho}
\bigl(z_4,\gamma(\mathcal H_1)\bigr)
\neq
\widehat\Omega_{\mathcal P,\rho}
\bigl(z_4,\gamma(\mathcal H_2)\bigr).
\]
Alors $z_4$ n'est pas un résumé suffisant pour cette famille et cette
résolution. L'expérience détecte l'activité d'une coordonnée manquante.
En revanche, le Théorème~\ref{thm:K-equivalence-plasticite} montre qu'aucune expérience fondée sur
les exécutions ne peut décider si cette coordonnée est décrite comme un
paramètre plastique de la loi ou comme une variable d'état.

\resultstep{cor:K-horloge-ontologie}
\begin{encadre}[title={Corollaire \theproposition{} --- Détectable comme horloge, indécidable comme ontologie}]
À $(\mathfrak B,\rho)$ fixés, l'observation peut établir :
\begin{enumerate}
\item l'insuffisance d'un résumé déclaré ;
\item l'activité décisionnelle d'une coordonnée supplémentaire ;
\item éventuellement sa loi d'usage et son échelle de temps.
\end{enumerate}
Elle ne peut pas établir, à partir des seules exécutions, si cette
coordonnée relève intrinsèquement de la loi ou de l'état.
\end{encadre}

\section{Fermeture du résumé étendu}
\label{annexe-K-fermeture}

Sous les hypothèses de l'Appendice I appliquées à l'état augmenté, il
existe une carte
\[
\widehat\Omega^{(5)}_{\mathcal P,\rho}:
Z_5\longrightarrow\mathcal Y_\rho,
\qquad
Z_5\subseteq\mathbb R^5,
\]
telle que
\[
\mathcal R^{\rho}_{z_5}(\mathcal P)
=
\widehat\Omega^{(5)}_{\mathcal P,\rho}(p,\chi,\theta,s,\gamma).
\]
Le résumé à quatre coordonnées est la projection
\[
\pi_4:Z_5\to Z_4,
\qquad
\pi_4(p,\chi,\theta,s,\gamma)=(p,\chi,\theta,s).
\]
Il est suffisant si et seulement si la carte de réponse se factorise par
$\pi_4$ sur le domaine considéré :
\[
\widehat\Omega^{(5)}_{\mathcal P,\rho}
=
\widehat\Omega^{(4)}_{\mathcal P,\rho}\circ\pi_4.
\]
Une paire témoin à même $z_4$ et issues différentes réfute exactement
cette factorisation.

\resultstep{prop:K-cinquieme-coordonnee}
\begin{encadre}[title={Proposition \theproposition{} --- Critère de nécessité de la cinquième coordonnée}]
Si deux états réalisables $z_5,z_5'$ vérifient
\[
\pi_4(z_5)=\pi_4(z_5')
\quad\text{et}\quad
\widehat\Omega^{(5)}_{\mathcal P,\rho}(z_5)\neq
\widehat\Omega^{(5)}_{\mathcal P,\rho}(z_5'),
\]
alors aucune carte définie seulement sur $Z_4$ ne peut reproduire la
réponse sur le domaine contenant ces deux états.
\end{encadre}

\paragraph{Démonstration.}
Si une telle carte existait, les deux états auraient la même image parce
qu'ils ont la même projection par $\pi_4$, en contradiction avec les
issues observées. \hfill$\square$

\section{Frontières de réponse sous présentation plastique}
\label{annexe-K-frontieres}

Pour une sortie binaire, posons
\[
A^{(5)}_{\mathcal P,\rho}
=
\{z_5\in Z_5:\widehat\Omega^{(5)}_{\mathcal P,\rho}(z_5)=1\},
\qquad
\Sigma^{(5)}_{\mathcal P,\rho}
=
\partial_{Z_5}A^{(5)}_{\mathcal P,\rho}.
\]
À $\gamma$ fixé, la section
\[
\Sigma_{\mathcal P,\rho}^{\gamma}
=
\Sigma^{(5)}_{\mathcal P,\rho}
\cap\{\gamma=\mathrm{constante}\}
\]
décrit la frontière vue par un système dont la coordonnée plastique est
maintenue à cette valeur pendant la sonde. Si la sonde modifie elle-même
$\gamma$, la réponse doit être calculée dans l'espace complet $Z_5$ : la
section figée n'est alors qu'un contrôle contrefactuel.

Deux histoires peuvent déposer le même $z_4$ mais des valeurs différentes
de $\gamma$. Elles interrogent alors des sections différentes de
$\Sigma^{(5)}$. Le déplacement résiduel de la frontière à $z_4$ égal est
une preuve d'activité de $\gamma$, non une preuve de sa localisation du
côté de la loi.

\resultstep{cor:K-invariance-frontiere}
\begin{encadre}[title={Corollaire \theproposition{} --- Invariance de la frontière complète}]
Les deux présentations du Théorème~\ref{thm:K-equivalence-plasticite} induisent le même ensemble
d'appropriation et la même frontière dans l'état augmenté $Z_5$. Ce qui
change est uniquement la manière de nommer la coordonnée $\gamma$.
\end{encadre}

\section{Décision de modélisation : frontière espèce/individu}
\label{annexe-K-espece-individu}

\begin{encadre}[title={D-08 --- Frontière déclarative}]
La distinction entre une loi d'espèce plastique et un individu doté
d'un état lent supplémentaire est relative à la présentation déclarée.
Cette convention est contrainte par l'équivalence dynamique : elle est un
choix de factorisation, jamais une licence pour modifier les exécutions.
Elle ne constitue pas un invariant de $\mathcal R^{\mathfrak B,\rho}$.
L'invariant accessible par le protocole est la structure des coordonnées
actives nécessaires à la fermeture du comportement sur le domaine et à
la résolution considérés.
\end{encadre}

D-08 ne supprime pas les conventions antérieures ; elle en précise la
portée. On peut continuer à dire qu'une histoire modifie la loi dans une
présentation donnée, à condition de ne pas transformer cette convention
en conclusion ontologique absolue.

\section{Spectre des horloges actives}
\label{annexe-K-spectre}

Le contenu positif laissé par le théorème est le \emph{spectre des
horloges actives}. À un niveau encore programmatique, il rassemble :
\begin{enumerate}
\item le nombre minimal de coordonnées nécessaires pour fermer la carte
de réponse sur le domaine déclaré ;
\item leurs échelles de temps caractéristiques ;
\item leurs lois de recharge, fatigue, oubli ou récupération ;
\item leur niveau d'activité ponctuelle, locale ou globale au sens de
l'Appendice I.
\end{enumerate}
Ce spectre est relatif à $(\mathfrak B,\rho)$ : une coordonnée inactive
pour une classe pauvre de sondes peut devenir nécessaire lorsque le
budget ou la résolution augmentent. La preuve de l'existence et de
l'unicité d'un spectre minimal dépasse la portée du présent appendice.

\section{Instance canonique à cinq horloges}
\label{annexe-K-instance}

L'extension canonique introduit un gain plastique $\gamma$ dans la loi
d'action, avec fatigue lors des événements chargés et récupération lente :
\[
\gamma^+=\gamma(1-r_\gamma),
\qquad
\dot\gamma=\varepsilon_\gamma(1-\gamma).
\]
La hiérarchie temporelle devient
\[
p<\chi<\theta<s<\gamma.
\]
Dans cette instance, deux histoires d'usage inégal peuvent être égalisées
en $(p,\chi,\theta,s)$ tout en laissant des valeurs distinctes de
$\gamma$. Une différence ultérieure de carte de sondes établit que
$z_4$ ne ferme plus le comportement, tandis que l'injection de
$z_5=(p,\chi,\theta,s,\gamma)$ reproduit la réponse.

Cette validation numérique est un témoin du Corollaire~\ref{cor:K-horloge-ontologie} et de la
Proposition~\ref{prop:K-cinquieme-coordonnee}. Elle n'entre pas dans la preuve du Théorème~\ref{thm:K-equivalence-plasticite},
laquelle repose uniquement sur l'identité des exécutions.

\section{Limites de portée}
\label{annexe-K-limites}

\paragraph{Mémoire non markovienne.}
Si la loi dépend d'une fonctionnelle de toute l'histoire qui ne possède
pas de réalisation d'état dans la classe admise, la construction par
état augmenté n'est pas disponible. Une réalisation de dimension infinie
peut parfois exister, mais elle change la catégorie d'espaces considérée.

\paragraph{Dépendance non causale.}
Une loi lisant une information future ne peut pas être remplacée par une
variable d'état causale sans modifier le protocole. Le théorème ne couvre
pas ce cas.

\paragraph{Domaines non réalisables.}
Une équivalence formelle définie sur des états inatteignables n'a pas de
portée empirique. Toutes les comparaisons doivent être restreintes au
domaine réalisable commun.

\paragraph{Bruits et lois de probabilité.}
Dans un système stochastique, l'identité trajectorielle exige un couplage
commun. Sans couplage, l'énoncé pertinent devient l'égalité des lois
d'exécution, puis des lois de réponse.

\paragraph{Observations dépendant de la présentation.}
Une lecture qui accéderait directement à l'étiquette interne ``paramètre
de loi'' violerait (P5). Elle ne serait plus une observation du système
par ses exécutions, mais une inspection du code ou de la convention de
modélisation.

\paragraph{Minimalité du résumé.}
Le théorème construit une réalisation, non une réalisation minimale. Plusieurs états augmentés de dimensions différentes peuvent être équivalents en réponse sur une classe de sondes donnée.

\paragraph{Changement de résolution.}
L'activité de $\gamma$ peut disparaître à résolution grossière et
réapparaître à résolution fine. Toutes les affirmations de nécessité et
de suffisance restent indexées par $(\mathfrak B,\rho)$.

\section{Dépendances logiques}
\label{annexe-K-dependances}

L'annexe utilise :
\begin{enumerate}
\item l'axiome d'observation et les équivalences de l'Appendice G ;
\item la réduction au contenu survivant de l'Appendice H ;
\item la fermeture du résumé lent et les niveaux d'activité de
l'Appendice I ;
\item la géométrie des frontières et la distinction pli/couture de
l'Appendice J.
\end{enumerate}
Elle prépare l'Appendice L : l'ultramétrique devra être formulée sur les
classes d'observation et rester invariante sous les rebaptêmes de
présentation établis ici.

\section{Registre de statut}
\label{annexe-K-registre}

\begin{description}
\item[Théorème~\ref{thm:K-equivalence-plasticite}.] Équivalence de présentation de la plasticité sous
réalisabilité causale et markovienne.
\item[Lemme~\ref{lem:K-identite-executions}.] Identité trajectorielle des deux présentations.
\item[Corollaire~\ref{cor:K-horloge-ontologie}.] Plasticité détectable comme horloge active et
indécidable comme ontologie.
\item[Proposition~\ref{prop:K-cinquieme-coordonnee}.] Critère de nécessité d'une coordonnée
supplémentaire.
\item[Corollaire~\ref{cor:K-invariance-frontiere}.] Invariance de l'ensemble d'appropriation et de la
frontière dans l'état augmenté.
\item[Décision D-08.] Frontière espèce/individu déclarative.
\item[Validation.] Instance canonique à cinq horloges ; elle illustre la
détection mais n'est pas utilisée dans la preuve.
\item[Question ouverte.] Caractériser les plasticités non réalisables par
un état causal markovien et définir un spectre minimal invariant des
horloges actives.
\end{description}

\section{Conclusion}
\label{annexe-K-conclusion}

La plasticité possède deux statuts qu'il faut maintenir simultanément.
Elle est observable lorsqu'elle rend insuffisant un résumé d'état et
modifie les réponses futures à état réduit égal. Mais sa localisation
entre loi et état n'est pas observable dès lors que les deux
présentations produisent les mêmes exécutions. Le résultat ne dissout pas
la plasticité : il déplace son contenu empirique vers ce qui résiste au
rebaptême, à savoir les coordonnées actives nécessaires, leurs échelles
de temps et leurs effets sur les frontières de réponse.
